% from ega3-6-fr.tex
\oldpage[III]{172}
Calculons les hypertor par la méthode de (6.6.2), en tenant compte de la remarque
(6.6.6); avec les notations de (6.6.2), chaque
$L_{\bullet\bullet}^{(i)}$ est un bicomplexe de $A_i$-modules, en désignant par
$A_i$ l'anneau de $Y^{(i)}$; il admet donc une résolution de Cartan--Eilenberg
projective $M_{\bullet\bullet\bullet}^{(i)}$ (au sens de (0, 11.7.1)) formée de
$A_i$-modules, et en vertu de (6.6.6), le premier membre de (6.10.1.1) est
canoniquement isomorphe à
\[
  \mathscr H_\bullet
  (M_{\bullet\bullet\bullet}\otimes_A Q'_\bullet),
\]
où
\[
  M_{\bullet\bullet\bullet}
  =
  M_{\bullet\bullet\bullet}^{(1)}
  \otimes_A M_{\bullet\bullet\bullet}^{(2)}
  \otimes_A\cdots\otimes_A M_{\bullet\bullet\bullet}^{(n)}
\]
et $\mathscr Q'_\bullet=\widetilde{Q'_\bullet}$. Comme, par hypothèse, les
anneaux $A_i$ sont des $A$-modules plats, les
$M_{\bullet\bullet\bullet}^{(i)}$ sont des tricomplexes de $A$-modules plats
($0_{\mathrm I}$, 6.2.1), et il en est de même de
$M_{\bullet\bullet\bullet}$; en outre, si $B$ est l'anneau de $Y$, produit
tensoriel des $A_i$, $M_{\bullet\bullet\bullet}$ est un tricomplexe de
$B$-modules; le complexe
\[
  \mathscr R_\bullet=(M_{\bullet\bullet\bullet})^\sim
\]
de $\mathscr O_Y$-Modules (où $M_{\bullet\bullet\bullet}$ est considéré comme
complexe simple) répond donc à la question, comme il résulte aisément de
(6.7.10).
\end{proof}

\begin{corollary}[6.10.2]
Dans l'énoncé de (6.10.1), on peut supposer $\mathscr R_\bullet$ limité
inférieurement. Lorsque les $\mathscr P_\bullet^{(i)}$ sont limités
supérieurement et les $Y^{(i)}$ de dimension cohomologique finie, on peut
supposer $\mathscr R_\bullet$ limité supérieurement.
\end{corollary}

\begin{proof}
La première assertion résulte de ce que les trois degrés de chacun des
$M_{\bullet\bullet\bullet}^{(i)}$ sont limités inférieurement; d'autre part, si
les anneaux $A_i$ sont de dimension cohomologique finie, le troisième degré de
chacun des $M_{\bullet\bullet\bullet}^{(i)}$ ne prend qu'un nombre fini de
valeurs, et il en est de même par construction de son premier degré (6.6.2);
comme son second degré est limité supérieurement s'il en est ainsi du degré de
$\mathscr P_\bullet^{(i)}$ (6.6.2), la seconde assertion en résulte aussitôt.
\end{proof}

\begin{remarks}[6.10.3]
\textnormal{(i)} Avec les notations de (6.10.1),
$\mathscr H_\bullet(\mathscr R_\bullet\otimes_S\mathscr Q'_\bullet)$ est
isomorphe à
$\mathbf{Tor}_\bullet(\mathscr R_\bullet,\mathscr Q'_\bullet)$ puisque
$\mathscr R_\bullet$ est formé de $\mathscr O_Y$-Modules $S$-plats (6.5.9); il
est donc (6.5.4) l'aboutissement d'une suite spectrale régulière de terme
$E_2$ donné par
\[
  \mathscr E_{pq}^2
  =
  \bigoplus_{q'+q''=q}
  \mathbf{Tor}_p^S
  \bigl(
    \mathscr H_{q'}(\mathscr R_\bullet),
    \mathscr H_{q''}(\mathscr Q'_\bullet)
  \bigr)
  \tag{6.10.3.1}\label{III.6.10.3.1.fr}
\]
qui n'est autre que la suite spectrale $(e)$ du changement de base (6.9.3).

\medskip
\textnormal{(ii)} Soit $\mathscr R'_\bullet$ un second complexe de
$\mathscr O_Y$-Modules quasi-cohérents, plat sur $S$, et soit
$g:\mathscr R'_\bullet\to\mathscr R_\bullet$ un homomorphisme de complexes tel
que
\[
  \mathscr H_\bullet(g):
  \mathscr H_\bullet(\mathscr R'_\bullet)
  \longrightarrow
  \mathscr H_\bullet(\mathscr R_\bullet)
\]
soit un isomorphisme. Alors, en vertu de (6.3.3) et (6.5.9), on déduit de $g$
un isomorphisme de $\partial$-foncteurs en $\mathscr Q'_\bullet$
\[
  \mathscr H_\bullet
  (\mathscr R'_\bullet\otimes_S\mathscr Q'_\bullet)
  \xrightarrow{\sim}
  \mathscr H_\bullet
  (\mathscr R_\bullet\otimes_S\mathscr Q'_\bullet)
\]
tel que le diagramme
\[
\xymatrix@C=28pt@R=30pt{
  \mathscr H_\bullet
  (\mathscr R'_\bullet\otimes_S\mathscr Q'_\bullet)
  \ar[r]^-{\sim}\ar[d]
  &
  \mathscr H_\bullet
  (\mathscr R_\bullet\otimes_S\mathscr Q'_\bullet)
  \ar[d]
  \\
  \mathscr H_\bullet
  (\mathscr R'_\bullet\otimes_Su^*(\mathscr Q'_\bullet))
  \ar[r]^-{\sim}
  &
  \mathscr H_\bullet
  (\mathscr R_\bullet\otimes_Su^*(\mathscr Q'_\bullet))
}
\]
\par\medskip\noindent
soit commutatif. Cela prouve donc que le complexe $\mathscr R_\bullet$ n'est
pas entièrement déterminé par les propriétés de (6.10.1).

\medskip
